\unitchapter{1}{10}{Higher Education System}{p1-10}

\unitmeta{Expected questions: \textasciitilde5 (10 marks) · Target: 4--5/5}

\begin{sylbox}

\begin{itemize}
\tightlist
\item[$\square$]
  Institutions of higher learning and education in ancient India
\item[$\square$]
  Evolution of higher learning and research in post-independence India
\item[$\square$]
  Oriental, conventional and non-conventional learning programmes in India
\item[$\square$]
  Professional, technical and skill-based education
\item[$\square$]
  Value education and environmental education
\item[$\square$]
  Policies, governance and administration
\end{itemize}

\end{sylbox}

\needspace{160pt}

\hypertarget{p1-10--1-ancient-indian-centres-of-learning}{%
\section{Ancient Indian Centres of Learning}\label{p1-10--1-ancient-indian-centres-of-learning}}

India had residential universities more than a thousand years before Europe. Questions usually ask who
founded a centre, where it was, which foreign scholars studied there, or who destroyed it.

\begingroup\small

\par\medskip\noindent\hfil\begin{tabular}{@{}
  >{\raggedright\arraybackslash}p{(\columnwidth - 4\tabcolsep) * \real{0.2126}}
  >{\raggedright\arraybackslash}p{(\columnwidth - 4\tabcolsep) * \real{0.2092}}
  >{\raggedright\arraybackslash}p{(\columnwidth - 4\tabcolsep) * \real{0.5782}}@{}}
\toprule\noalign{}
\begin{minipage}[b]{\linewidth}\raggedright
\textbf{Centre}
\end{minipage} & \begin{minipage}[b]{\linewidth}\raggedright
\textbf{Location (present)}
\end{minipage} & \begin{minipage}[b]{\linewidth}\raggedright
\textbf{Key facts}
\end{minipage} \\
\midrule\noalign{}
\textbf{Takshashila} & Rawalpindi, Pakistan & Oldest (\textasciitilde6th c. BCE); Chanakya, Panini, Charaka, Jivaka; no formal degrees \\
\textbf{Nalanda} & Bihar & Founded by \textbf{Kumaragupta I} (Gupta, 5th c. CE); Xuanzang \& I-tsing studied; library \emph{Dharmaganja} (Ratnasagara, Ratnodadhi, Ratnaranjaka); destroyed by \textbf{Bakhtiyar Khalji} (\textasciitilde1193) \\
\textbf{Vikramashila} & Bhagalpur, Bihar & Founded by \textbf{Dharmapala} (Pala, 8th c.); Tantric Buddhism; Atisha Dipankara \\
\textbf{Valabhi} & Gujarat & Maitraka dynasty; Hinayana Buddhism \\
\textbf{Odantapuri} & Bihar & Gopala (Pala) \\
\textbf{Somapura Mahavihara} & Bangladesh & Dharmapala; UNESCO site \\
\textbf{Kanchi (Kanchipuram)} & Tamil Nadu & Pallavas; Ghatika \\
\textbf{Pushpagiri} & Odisha & \\
\textbf{Jagaddala} & Bengal & Ramapala \\
\bottomrule\noalign{}
\end{tabular}\hfil\null\par\medskip


\endgroup

Vedic education: \textbf{Gurukul}; teaching via \emph{Shravana (listening) → Manana (reflection) → Nididhyasana (meditation)}. Upanayana ceremony marked entry.
Buddhist education: \emph{Pabbajja} (entry) and \emph{Upasampada} (full ordination) ceremonies; viharas \& monasteries.

\needspace{125pt}

\hypertarget{p1-10--2-evolution-of-higher-education--timeline}{%
\section{Evolution of Higher Education --- Timeline}\label{p1-10--2-evolution-of-higher-education--timeline}}

Modern higher education in India has been shaped by a succession of commissions and policies. Their order
and their key recommendations are among the most frequently asked facts in this unit.

\begin{figure}[!htbp]
\centering
\begin{tikzpicture}[y=-0.6cm]
  \draw[line width=1.6pt, fslate] (0,0.3) -- (0,19.7);
  \foreach \y/\yr/\ev in {1/1781/{Calcutta Madrasa (Warren Hastings)}, 2/1791/{Sanskrit College, Benares (Jonathan Duncan)},
      3/1813/{Charter Act — ₹1 lakh a year for education}, 4/1835/{Macaulay's Minute — English education},
      5/1854/{Wood's Despatch — "Magna Carta of English education in India"}, 6/1857/{Universities of Calcutta, Bombay and Madras},
      7/1882/{Hunter Commission}, 8/1902/{Raleigh Commission → Indian Universities Act, 1904},
      9/1917/{Sadler (Calcutta University) Commission}, 10/1948/{Radhakrishnan Commission — recommended the UGC},
      11/1953/{UGC set up (inaugurated 28 December 1953)}, 12/1956/{UGC Act — UGC becomes a statutory body},
      13/1964/{Kothari Commission (1964–66) — 10+2+3 pattern}, 14/1968/{First National Policy on Education},
      15/1986/{National Policy on Education (revised 1992, Programme of Action)}, 16/2005/{National Knowledge Commission (Sam Pitroda)},
      17/2009/{Yash Pal Committee}, 18/2013/{RUSA — funding for state institutions}, 19/2020/{National Education Policy 2020 (K. Kasturirangan)}} {
    \filldraw[fill=fwine, draw=white, line width=0.8pt] (0,\y) circle (3.2pt);
    \node[anchor=east, font=\small\bfseries] at (-0.3,\y) {\yr};
    \node[anchor=west, font=\small, text width=125mm] at (0.3,\y) {\ev};
  }
\end{tikzpicture}
\caption{The evolution of modern higher education in India.}
\end{figure}


\needspace{125pt}

\hypertarget{p1-10--3-nep-2020--higher-education-highlights-most-asked}{%
\section{NEP 2020 --- Higher Education Highlights (most asked)}\label{p1-10--3-nep-2020--higher-education-highlights-most-asked}}

The National Education Policy 2020 is the first comprehensive policy since 1986. Its higher-education
reforms aim at flexibility for students, multidisciplinary institutions and lighter but tighter regulation.

\begin{figure}[!htbp]
\centering
\begin{tikzpicture}
  \node[circ, minimum size=22mm, fill=fwine!15, font=\small\bfseries, align=center] (N) at (0,0) {NEP\\2020};
  \node[slate, text width=40mm, align=left, font=\footnotesize] (a) at (-4.6,1.7) {\textbf{Structure}\\5+3+3+4 schooling\\GER 50\% by 2035\\multidisciplinary institutions};
  \node[sage, text width=40mm, align=left, font=\footnotesize] (b) at (4.6,1.7) {\textbf{Degrees}\\4-year UG, multiple entry–exit\\certificate–diploma–degree–honours\\M.Phil discontinued};
  \node[sand, text width=40mm, align=left, font=\footnotesize] (c) at (-4.6,-1.7) {\textbf{Institutions}\\research \& teaching universities\\autonomous degree-granting colleges\\affiliation phased out by 2035};
  \node[rose, text width=40mm, align=left, font=\footnotesize] (d) at (4.6,-1.7) {\textbf{Others}\\Academic Bank of Credits\\National Research Foundation\\6\% of GDP on education};
  \foreach \n in {a,b,c,d} \draw[thinline] (N) -- (\n);
  \node[grey, minimum width=34mm] (H) at (0,-4.1) {\textbf{HECI} — single regulator};
  \foreach \x/\t/\s in {-4.8/NHERC/regulation, -1.6/NAC/accreditation, 1.6/HEGC/funding, 4.8/GEC/standards}
    \node[plain, minimum width=26mm, font=\footnotesize] (v\t) at (\x,-5.5) {\textbf{\t}\\\s};
  \foreach \t in {NHERC,NAC,HEGC,GEC} \draw[arr] (H.south) -- (v\t.north);
  \draw[thinline] (N) -- (H);
\end{tikzpicture}
\caption{The main reforms of NEP 2020 for higher education, and the four verticals of the proposed Higher
Education Commission of India.}
\end{figure}


\begin{itemize}
\tightlist
\item
  \textbf{HECI} (Higher Education Commission of India) --- single umbrella regulator with 4 verticals: NHERC, NAC, HEGC, GEC (proposed; would subsume UGC/AICTE/NCTE except medical \& legal).
\item
  \textbf{ANRF} (Anusandhan National Research Foundation) Act 2023 --- implements NRF idea.
\item
  \textbf{Multidisciplinary Education \& Research Universities (MERUs)}.
\end{itemize}

\needspace{130pt}

\hypertarget{p1-10--4-types-of-institutions--programmes}{%
\section{Types of Institutions \& Programmes}\label{p1-10--4-types-of-institutions--programmes}}

\begingroup\small

\par\medskip\noindent\hfil\begin{tabular}{@{}
  >{\raggedright\arraybackslash}p{(\columnwidth - 2\tabcolsep) * \real{0.3261}}
  >{\raggedright\arraybackslash}p{(\columnwidth - 2\tabcolsep) * \real{0.6739}}@{}}
\toprule\noalign{}
\begin{minipage}[b]{\linewidth}\raggedright
\textbf{Type}
\end{minipage} & \begin{minipage}[b]{\linewidth}\raggedright
\textbf{Notes}
\end{minipage} \\
\midrule\noalign{}
Central university & Act of Parliament (e.g. JNU, DU, BHU, AMU) \\
State university & State Act \\
Deemed-to-be university & Sec 3 of UGC Act, declared by Central Govt on UGC advice \\
Private university & State Act; cannot affiliate colleges \\
Institutions of National Importance & Act of Parliament (IITs, NITs, AIIMS, IIMs) \\
Institution of Eminence & 2017 scheme --- 10 public + 10 private \\
Open universities & IGNOU (1985); first state open univ: Dr. B.R. Ambedkar Open Univ, Hyderabad (1982) \\
\bottomrule\noalign{}
\end{tabular}\hfil\null\par\medskip


\endgroup

\begin{itemize}
\tightlist
\item
  \textbf{Oriental learning}: Sanskrit, Pali, Persian, Arabic institutions --- e.g. Rashtriya Sanskrit Sansthan; Central Sanskrit University (2020).
\item
  \textbf{Conventional}: classroom-based degrees. \textbf{Non-conventional}: distance, open, online, MOOCs, part-time.
\end{itemize}

\needspace{160pt}

\hypertarget{p1-10--5-regulatory--professional-bodies}{%
\section{Regulatory \& Professional Bodies}\label{p1-10--5-regulatory--professional-bodies}}

India regulates higher education through a separate body for each sector. Matching each body with its
field and its year of establishment is a standard question.

\begingroup\small

\par\medskip\noindent\hfil\begin{tabular}{@{}
  >{\raggedright\arraybackslash}p{(\columnwidth - 4\tabcolsep) * \real{0.2006}}
  >{\raggedright\arraybackslash}p{(\columnwidth - 4\tabcolsep) * \real{0.5362}}
  >{\raggedright\arraybackslash}p{(\columnwidth - 4\tabcolsep) * \real{0.2161}}@{}}
\toprule\noalign{}
\begin{minipage}[b]{\linewidth}\raggedright
\textbf{Body}
\end{minipage} & \begin{minipage}[b]{\linewidth}\raggedright
\textbf{Field}
\end{minipage} & \begin{minipage}[b]{\linewidth}\raggedright
\textbf{Est.}
\end{minipage} \\
\midrule\noalign{}
UGC & Universities (funding, standards) & 1953 / Act 1956 \\
AICTE & Technical education & 1945 (statutory 1987) \\
NCTE & Teacher education & 1973 (statutory 1993) \\
NMC & Medical (replaced MCI) & 2020 \\
BCI & Legal education & 1961 \\
ICAR & Agriculture & 1929 \\
PCI / INC / DCI & Pharmacy / Nursing / Dental & \\
COA & Architecture & 1972 \\
RCI & Rehabilitation & 1992 \\
\textbf{NAAC} & Accreditation of HEIs (Bengaluru) & 1994 \\
\textbf{NBA} & Accreditation of programmes (technical) & 1994 \\
NIRF & Ranking (MoE) & 2015 \\
NTA & Testing (UGC NET, CUET, NEET, JEE Main) & 2017 \\
CABE & Central Advisory Board of Education --- oldest advisory body & 1920/1935 \\
NIEPA & Educational planning \& admin & \\
\bottomrule\noalign{}
\end{tabular}\hfil\null\par\medskip


\endgroup

NAAC grades: A++, A+, A, B++, B+, B, C, D on 4-point CGPA; criteria = \textbf{7} (Curricular aspects; Teaching-learning \& evaluation; Research, innovations \& extension; Infrastructure; Student support; Governance; Institutional values \& best practices).

\needspace{95pt}

\hypertarget{p1-10--6-skill-based--professional-education}{%
\section{Skill-based \& Professional Education}\label{p1-10--6-skill-based--professional-education}}

\begin{itemize}
\tightlist
\item
  \textbf{Skill India / PMKVY} (2015), \textbf{NSQF} (National Skills Qualification Framework, levels 1--10), \textbf{NCVET} (2018), B.Voc / D.Voc, Community colleges, \textbf{Apprenticeship embedded degree programmes}.
\item
  \textbf{NHEQF} (National Higher Education Qualification Framework, levels 4.5--8) \& \textbf{NCrF} (National Credit Framework).
\end{itemize}

\needspace{95pt}

\hypertarget{p1-10--7-value--environmental-education}{%
\section{Value \& Environmental Education}\label{p1-10--7-value--environmental-education}}

\begin{itemize}
\tightlist
\item
  Value education: truth, non-violence, peace, love, right conduct (Sathya Sai model); constitutional values; NEP emphasises ethics, Indian Knowledge Systems (IKS).
\item
  Environmental education made compulsory at all levels by \textbf{Supreme Court (M.C. Mehta case, 1991/2003)}; UGC mandated a 6-month core module course for UG.
\item
  \textbf{Tbilisi Declaration (1977)} --- first intergovernmental conference on environmental education (UNESCO-UNEP); Belgrade Charter (1975).
\end{itemize}

\needspace{125pt}

\hypertarget{p1-10--8-governance--administration}{%
\section{Governance \& Administration}\label{p1-10--8-governance--administration}}

A university is governed by a hierarchy of officers and statutory bodies defined in its Act. The figure
shows the usual arrangement; details vary between central and state universities.

\begin{figure}[!htbp]
\centering
\begin{adjustbox}{max width=\linewidth}
\begin{tikzpicture}
  \node[slate, text width=50mm] (V) at (0,0) {\textbf{Visitor} — President of India\\\footnotesize central universities};
  \node[slate, text width=50mm] (C) at (0,-1.4) {\textbf{Chancellor} — Governor\\\footnotesize state universities};
  \node[sand, text width=50mm] (VC) at (0,-2.8) {\textbf{Vice-Chancellor}\\\footnotesize chief executive and academic officer};
  \draw[arr] (V) -- (C); \draw[arr] (C) -- (VC);
  \foreach \x/\t/\s [count=\k] in {-6.0/{Executive\\Council}/{apex executive}, -3.6/{Academic\\Council}/{apex academic}, -1.2/{Court /\\Senate}/{}, 1.2/Registrar/{administration}, 3.6/{Finance\\Officer}/{}, 6.0/{Controller of\\Examinations}/{}} {
    \node[sage, text width=20mm, font=\footnotesize, anchor=north] (b\k) at (\x,-4.2) {\textbf{\t}\\\s};
    \draw[arr] (VC.south) -- (b\k.north);
  }
\end{tikzpicture}
\end{adjustbox}
\caption{The governance structure of an Indian university.}
\end{figure}


\begin{itemize}
\tightlist
\item
  Education was moved from State List to \textbf{Concurrent List} by the \textbf{42nd Amendment (1976)}.
\item
  \textbf{Article 21A} --- right to education (6--14 yrs) via 86th Amendment 2002; RTE Act 2009.
\item
  Article 30 --- minorities\textquotesingle{} right to establish institutions; Article 45 --- ECCE; Article 46 --- education of SC/ST/weaker sections.
\item
  Entry 66, Union List --- coordination \& determination of standards in higher education.
\item
  \textbf{Ministry of Education} (renamed from MHRD in 2020). \textbf{RUSA (2013)} --- funding to state HEIs.
\end{itemize}

\needspace{131pt}

\hypertarget{p1-10--9-deeper-dive--nep-2020-details-quality-frameworks--key-regulations}{%
\section{Deeper Dive --- NEP 2020 Details, Quality Frameworks \& Key Regulations}\label{p1-10--9-deeper-dive--nep-2020-details-quality-frameworks--key-regulations}}

\needspace{95pt}

\hypertarget{p1-10--91-nep-2020--multiple-entryexit--credit-system}{%
\subsection{NEP 2020 --- Multiple Entry/Exit \& Credit System}\label{p1-10--91-nep-2020--multiple-entryexit--credit-system}}

\begin{figure}[!htbp]
\centering
\begin{tikzpicture}
  \foreach \i/\t/\c/\col in {0/{UG certificate}/40/fslate, 1/{UG diploma}/80/fsteel, 2/{Bachelor's degree}/120/fsage, 3/{Honours, or honours with research}/160/fsand} {
    \pic at ({\i*2.9},{\i*0.75}) {slab={2.9}{0.75}{0.9}{\col}};
    \node[font=\footnotesize\bfseries] at ({\i*2.9+1.45},{\i*0.75+0.375}) {Year \pgfmathparse{int(\i+1)}\pgfmathresult};
    \node[font=\footnotesize, text width=26mm, align=center, anchor=north] at ({\i*2.9+1.45},-0.15) {exit: \t\\\textit{\c\ credits}};
  }
  \node[rose, text width=27mm, font=\footnotesize] (pg2) at (14.3,2.75) {\textbf{2-year master's}\\after 3 years};
  \node[rose, text width=27mm, font=\footnotesize] (pg1) at (14.3,4.5) {\textbf{1-year master's}\\after 4 years};
  \node[rose, text width=27mm, font=\footnotesize] (phd) at (14.3,0.9) {\textbf{PhD directly}\\honours with research, CGPA $\ge 7.5$};
  \draw[arr] (11.9,3.55) |- (pg1.west);
  \draw[arr] (9.0,2.85) -- (pg2.west);
  \draw[arr] (12.2,2.2) |- (phd.west);
\end{tikzpicture}
\caption{NEP 2020's four-year undergraduate programme, with an exit after every year and the routes onward.}
\end{figure}


\begin{itemize}
\tightlist
\item
  \textbf{UGC Curriculum \& Credit Framework for UG Programmes (CCFUP, 2022)}; credits stored in \textbf{Academic Bank of Credits (ABC)}, ID via DigiLocker / APAAR ID (One Nation One Student ID).
\item
  Students may earn up to \textbf{40\% of credits} through SWAYAM/online courses per semester (UGC Credit Framework for Online Learning Courses through SWAYAM, 2021).
\item
  Course categories: Major (core), Minor, Multidisciplinary, Ability Enhancement (AEC), Skill Enhancement (SEC), Value-Added (VAC), internships, research project.
\item
  \textbf{UGC (Minimum Standards and Procedures for Award of PhD) Regulations 2022}: 4-year UG (with 7.5 CGPA) eligible for PhD; M.Phil not required; 2-year coursework rule relaxed; publication requirement removed.
\item
  \textbf{Professor of Practice} scheme (2022): industry experts can teach (up to 10\% of sanctioned posts).
\item
  Foreign university campuses permitted (UGC Regulations 2023); GIFT City.
\item
  \textbf{Light but tight} regulation; graded autonomy; Board of Governors for institutions.
\end{itemize}

\needspace{130pt}

\hypertarget{p1-10--92-nirf--parameters-since-2016}{%
\subsection{NIRF --- Parameters (since 2016)}\label{p1-10--92-nirf--parameters-since-2016}}

\begingroup\small

\par\medskip\noindent\hfil\begin{tabular}{@{}
  >{\raggedright\arraybackslash}p{(\columnwidth - 2\tabcolsep) * \real{0.3556}}
  >{\raggedright\arraybackslash}p{(\columnwidth - 2\tabcolsep) * \real{0.2531}}@{}}
\toprule\noalign{}
\begin{minipage}[b]{\linewidth}\raggedright
\textbf{Parameter}
\end{minipage} & \begin{minipage}[b]{\linewidth}\raggedright
\textbf{Weight (overall category)}
\end{minipage} \\
\midrule\noalign{}
Teaching, Learning \& Resources (TLR) & 30\% \\
Research \& Professional Practice (RP) & 30\% \\
Graduation Outcomes (GO) & 20\% \\
Outreach \& Inclusivity (OI) & 10\% \\
Perception (PR) & 10\% \\
\bottomrule\noalign{}
\end{tabular}\hfil\null\par\medskip


\endgroup

\needspace{130pt}

\hypertarget{p1-10--93-naac-grading-cumulative-gpa-on-a-4-point-scale}{%
\subsection{NAAC Grading (cumulative GPA on a 4-point scale)}\label{p1-10--93-naac-grading-cumulative-gpa-on-a-4-point-scale}}

\begingroup\small

\par\medskip\noindent\hfil\begin{tabular}{@{}
  >{\raggedright\arraybackslash}p{(\columnwidth - 4\tabcolsep) * \real{0.1434}}
  >{\raggedright\arraybackslash}p{(\columnwidth - 4\tabcolsep) * \real{0.0739}}
  >{\raggedright\arraybackslash}p{(\columnwidth - 4\tabcolsep) * \real{0.1425}}@{}}
\toprule\noalign{}
\begin{minipage}[b]{\linewidth}\raggedright
\textbf{CGPA range}
\end{minipage} & \begin{minipage}[b]{\linewidth}\raggedright
\textbf{Grade}
\end{minipage} & \begin{minipage}[b]{\linewidth}\raggedright
\textbf{Status}
\end{minipage} \\
\midrule\noalign{}
3.51--4.00 & A++ & Accredited \\
3.26--3.50 & A+ & Accredited \\
3.01--3.25 & A & Accredited \\
2.76--3.00 & B++ & Accredited \\
2.51--2.75 & B+ & Accredited \\
2.01--2.50 & B & Accredited \\
1.51--2.00 & C & Accredited \\
≤ 1.50 & D & Not accredited \\
\bottomrule\noalign{}
\end{tabular}\hfil\null\par\medskip


\endgroup

NAAC is moving to \textbf{binary accreditation} (accredited / not accredited) and \textbf{Maturity-Based Graded Levels (1--5)} following the Radhakrishnan Committee (2022--24) recommendations.

\needspace{130pt}

\hypertarget{p1-10--94-constitutional--legal-provisions-on-education}{%
\subsection{Constitutional \& Legal Provisions on Education}\label{p1-10--94-constitutional--legal-provisions-on-education}}

\begingroup\small

\par\medskip\noindent\hfil\begin{tabular}{@{}
  >{\raggedright\arraybackslash}p{(\columnwidth - 2\tabcolsep) * \real{0.2267}}
  >{\raggedright\arraybackslash}p{(\columnwidth - 2\tabcolsep) * \real{0.5811}}@{}}
\toprule\noalign{}
\begin{minipage}[b]{\linewidth}\raggedright
\textbf{Provision}
\end{minipage} & \begin{minipage}[b]{\linewidth}\raggedright
\textbf{Content}
\end{minipage} \\
\midrule\noalign{}
Article 21A & Free \& compulsory education 6--14 years (86th Amendment, 2002) \\
Article 28 & No religious instruction in fully state-funded institutions \\
Article 29 & Protection of minorities\textquotesingle{} cultural \& educational interests \\
Article 30 & Minorities\textquotesingle{} right to establish \& administer institutions \\
Article 45 & Early childhood care \& education below 6 years (DPSP) \\
Article 46 & Promote education of SC/ST \& weaker sections \\
Article 51A(k) & Parent\textquotesingle s duty to provide education to child 6--14 \\
Article 350A & Instruction in mother tongue at primary stage \\
Union List entry 63--66 & Institutions of national importance; standards in higher education \\
Concurrent List entry 25 & Education (since 42nd Amendment) \\
\bottomrule\noalign{}
\end{tabular}\hfil\null\par\medskip


\endgroup

\needspace{95pt}

\hypertarget{p1-10--95-value-education--indian-knowledge-systems}{%
\subsection{Value Education \& Indian Knowledge Systems}\label{p1-10--95-value-education--indian-knowledge-systems}}

\begin{itemize}
\tightlist
\item
  \textbf{IKS division} at AICTE (2020); Bharatiya Bhasha Samiti; courses on Indian philosophy, mathematics, astronomy, Ayurveda.
\item
  Value frameworks: Delors Commission (UNESCO, 1996) "\textbf{Learning: The Treasure Within}" --- \textbf{four pillars}: learning to know, to do, to live together, to be.
\item
  \textbf{Mulya Pravah} (UGC, 2019, revised 2.0 in 2023): guidelines for human values \& professional ethics in HEIs.
\item
  \textbf{Deeksharambh}: student induction programme (UGC).
\item
  \textbf{Paramarsh} (mentoring for NAAC accreditation), \textbf{STRIDE} (research), \textbf{LOCF} (learning outcomes-based curriculum framework).
\end{itemize}

\needspace{95pt}

\hypertarget{p1-10--previous-year-questions-pyq-pattern}{%
\section{Previous Year Questions (PYQ pattern)}\label{p1-10--previous-year-questions-pyq-pattern}}

\begin{enumerate}
\def\labelenumi{\arabic{enumi}.}
\tightlist
\item
  Nalanda University was founded by: \textbf{Kumaragupta I}
\item
  Vikramashila was founded by: \textbf{Dharmapala}
\item
  Which Chinese traveller studied at Nalanda? \textbf{Xuanzang (Hiuen Tsang)} and I-tsing
\item
  Nalanda was destroyed by: \textbf{Bakhtiyar Khalji}
\item
  Takshashila is located in present-day: \textbf{Pakistan}
\item
  "Magna Carta of English education in India": \textbf{Wood\textquotesingle s Despatch 1854}
\item
  The first three universities were established in: \textbf{1857}
\item
  University Education Commission (1948) was chaired by: \textbf{Dr. S. Radhakrishnan}
\item
  10+2+3 pattern was recommended by: \textbf{Kothari Commission}
\item
  UGC became a statutory body in: \textbf{1956}
\item
  NEP 2020 GER target for higher education: \textbf{50\% by 2035}
\item
  Which degree is discontinued by NEP 2020? \textbf{M.Phil}
\item
  NEP 2020 proposes a single regulator: \textbf{HECI}
\item
  NAAC is located in: \textbf{Bengaluru}
\item
  Programme-level accreditation for engineering: \textbf{NBA}
\item
  Education was placed in the Concurrent List by: \textbf{42nd Amendment, 1976}
\item
  Right to Education is under Article: \textbf{21A}
\item
  IGNOU was established in: \textbf{1985}
\item
  Apex academic body of a university: \textbf{Academic Council}
\item
  Who is the Visitor of central universities? \textbf{President of India}
\item
  Number of NAAC criteria: \textbf{7}
\item
  The Tbilisi conference (1977) dealt with: \textbf{Environmental education}
\item
  Arrange chronologically: Macaulay\textquotesingle s Minute (1835), Wood\textquotesingle s Despatch (1854), Hunter (1882), Sadler (1917), Radhakrishnan (1948), Kothari (1964--66)
\end{enumerate}

\textbf{More practice questions}

\begin{enumerate}
\def\labelenumi{\arabic{enumi}.}
\setcounter{enumi}{23}
\tightlist
\item
  Credits required for a UG certificate under NEP\textquotesingle s multiple exit: \textbf{40}
\item
  Maximum credits that can be earned through SWAYAM per semester (UGC 2021): \textbf{40\%}
\item
  "Learning: The Treasure Within" is the report of the: \textbf{Delors Commission (UNESCO, 1996)}
\item
  Four pillars of learning do NOT include: (a) to know (b) to do (c) to compete (d) to be --- \textbf{Ans: (c)}
\item
  Weight of Teaching, Learning \& Resources in NIRF overall ranking: \textbf{30\%}
\item
  NAAC grade for CGPA 3.30: \textbf{A+}
\item
  Article 30 deals with: \textbf{right of minorities to establish educational institutions}
\item
  Article 350A provides for: \textbf{instruction in mother tongue at the primary stage}
\item
  Mulya Pravah guidelines relate to: \textbf{human values and professional ethics in HEIs}
\item
  Student induction programme of UGC: \textbf{Deeksharambh}
\item
  UGC scheme for mentoring institutions towards accreditation: \textbf{Paramarsh}
\item
  Under UGC PhD Regulations 2022, a 4-year UG graduate needs a minimum CGPA of: \textbf{7.5}
\end{enumerate}

\begin{revbox}

\begin{itemize}
\tightlist
\item
  Nalanda--Kumaragupta I · Vikramashila--Dharmapala · Destroyer--Khalji
\item
  1835 Macaulay · 1854 Wood · 1857 3 univs · 1948 Radhakrishnan · 1953/56 UGC · 1964-66 Kothari · 1968/1986 NPE · 2020 NEP
\item
  NEP: GER 50\% (2035), 4-yr UG, ABC, HECI (4 verticals), M.Phil out
\end{itemize}

\end{revbox}
